% ===== AI-DRAFTED (pending author review) =====
\section{Proportionality, the Core, and Stability}\label{sec:prop}
\aibanner

The notions above are imposed by a designer who fixes demographic constraints in advance. A complementary, \emph{game-theoretic} family of definitions instead lets \emph{coalitions of individuals} certify their own fair treatment, removing the need to pre-specify protected groups.

\subsection{Proportional fairness and the core}
\citet{chen2019proportionally} introduced \emph{proportional fairness} for centroid clustering, borrowing the notion of the \emph{core} from cooperative game theory: any sufficiently large group of points is entitled to its own center. Given $k$ centers, a subset $S\subseteq P$ with $|S|\ge \lceil n/k\rceil$ is a \emph{blocking coalition} if there exists a point $y\in P$ such that $d(x,y)<d(x,C)$ for every $x\in S$---all members of $S$ would strictly prefer to be served by $y$. A solution is \emph{proportionally fair} (lies in the core) if it admits no blocking coalition.

Exact proportional solutions need not exist, which motivates $\rho$-approximate proportionality, where a coalition is blocking only if every member improves by more than a factor $\rho$. \citet{chen2019proportionally} showed that a simple \emph{greedy capture} procedure achieves a $(1+\sqrt{2})$-approximate core for the $k$-median-style objective, and proved separations from classical optimality. \citet{DBLP:conf/icalp/Micha020} revisited the model, tightening bounds for Euclidean instances and clarifying the price of proportionality relative to the optimal $k$-median cost.

Because proportional fairness is defined without reference to protected attributes, it is attractive when group labels are unavailable or contested; the trade-off is that the guarantee concerns \emph{coalitions} rather than demographic balance, and the two need not coincide. A active recent line sharpens this picture: improved lower bounds on the achievable core approximation~\citep{proplowerbounds2026}, connections between proportional clustering, the $\beta$-plurality problem, and metric distortion~\citep{propplurality2025}, a \emph{proportionally representative} variant that constrains the chosen centers rather than the assignment~\citep{proprepclustering2023}, and a framework unifying proportional fairness across centroid and non-centroid (e.g.\ correlation-style) objectives~\citep{unifyingprop2026}.

\subsection{Preference- and stability-based fairness}
A related line treats each individual's \emph{preference} over clusters as the fairness primitive. \citet{ahmadi2022individual} formalized \emph{individual-preference (IP) stability}: a clustering is IP-stable if no point's average distance to its own cluster exceeds its average distance to any other cluster, so that no individual would, on average, prefer to defect. They characterize when such clusterings exist and give approximation algorithms---exact on the line and bounded-factor in general metrics---together with hardness results, with subsequent work giving \emph{scalable} algorithms for IP-stable clustering~\citep{ipstablescalable2024}. Conceptually, these definitions import \emph{envy-freeness} and stable matching into clustering, linking fair clustering to the broader literature on stability in coalition formation.
